\chapter{Methodology}
\label{methodology}

The present work consists of a translation and PE experiment from English into German with 12 professional translators. The selected method combines non-participant observation, screen recordings, retrospective interviews, and target text annotation. It was inspired by the field of TPR and \textcite{albl2017professional}, who combine screen recordings and retrospective interviews in research on the translation of non-native English in the context of the European Parliament. TPR focuses on the translation process and typical research questions include problem-solving, decision-making, revision, time pressure, and motivation \citep[123f]{jaaskelainen2011studying}. It borrows methodologies from other disciplines, such as psychology, linguistics, and the cognitive sciences (\cite[18]{alves2015translation}). Usually, data on the translation process are collected through keylogging, eye tracking and/or screen recordings, and triangulated with verbal reports, such as think-aloud protocols (TAP) or retrospective interviews. In the last few years, new technologies such as electroencephalograms (EEGs) have also been applied in research settings \citep[21]{jakobsen2017translation}.

The experiment was conducted between July and October 2022. Twelve professional translators were recruited (see more details on the study participants in Subsection \ref{sec:participants}). Before the experiment, participants received instructions on the tasks, PE guidelines by the Translation Automation User Society (TAUS), \footnote{\url{https://info.taus.net/mt-PE-guidelines}} and a handout on various strategies for GFL to prepare for their participation. Participants also completed an online questionnaire to collect demographic data as well as data on their use of GFL and its perceived difficulty.

Participants were divided into two groups. The first group was asked to translate three brief texts of approximately 150 words each on three English-language TV series that feature non-binary actors and characters (see Subsection \ref{sec:assignments}). The second group was asked to manually post-edit the same texts that had been previously machine-translated, i.e. participants received a text document divided into two columns, one containing the English language source text and the other containing the MT. The three texts were retrieved online, and in the case of PE, translated with the freely available MT system DeepL\footnote{\url{https://www.deepl.com/translator}} in June 2022. DeepL was selected for this study because it is a widely known and used commercial system. For each text, participants were instructed to adopt a different approach to GFL, that is, (i) gender-neutral rewording, (ii) gender-inclusive characters, and  (iii) neosystems. For each approach, they could freely select specific strategies from the provided handout, e.g. gender star (*) or underscore (\_) amongst others for (ii). The handout as well as other electronic resources such as online dictionaries were allowed during the experiment. 

The study was conducted online for several reasons. It aimed for an authentic and nonintrusive experimental setting in which translators could work from a familiar environment, notably their office at home, with familiar computers and, if desired, known CAT tools. The only restriction was to use one screen for the whole process to be recorded, and the translating participants could not use MT. The target country for participant recruitment was Austria because UNIVERSITAS, the Austrian Interpreters and Translators Association, financially supported the study. Due to difficulties in finding a large number of participants, especially non-binary people, who are still a marginalised group, recruitment was not restricted to one specific city. These difficulties are common to research on non-binary GFL (see \cite{burtscher-2022-genderfairmt,gromann2023participatory}), probably because it is still a controversial, political, and ideologically loaded topic.
%Whereas in controlled settings much preciser methodologies such as key-logging and eye tracking can be used to gain insights in the translations process (see for example \cite{alves2009new}), they present other limitations and reduce ecological validity (see for example \cite{o2009eye}). Finally, the online setting was also preferred because changing covid-regulations as well as delay to due participants' infection, that would have complicated the constant access to a laboratory, were feared. Nevertheless, these and other limitations of the present study are further discussed in Chapter \ref{sec:limitations}. 

Participants joined a video conference and were asked to share their screens as in Figure \ref{fig:screenshot_experiment}. During the translation and PE process, the video conference was open in the background and the shared screen was recorded. Participants could decide whether and when to make pauses and did not have to work under time constraints. After the translation and PE tasks, they were interviewed on their impressions, strategies, and challenging aspects of the study, amongst other things. They also discussed the possibility of using MT to translate GFL.

\begin{figure}[ht]
    \centering
    \includegraphics[width=1.0\textwidth]{figures/screenshot_experiment.pdf}
    \caption{Study participant's screen during the translation process}
    \label{fig:screenshot_experiment}
\end{figure}

The analysis of the translation and PE process was inspired by Krings' (\cite*{krings2001repairing}) division into temporal, technical, and cognitive effort. In order to shed light on the cognitive processes involved in integrating GFL in translation and PE, it focused on the translation and PE times, the screen activity tracking data, and participants' experiences and impressions. These different aspects of the translation and PE process allowed investigation of the effort as well as strengths and limitations of the three GFL strategies tested in the present work.

Screen recordings were used to measure translation and PE times. Though they are usually combined with other methods, such as keylogging (e.g. \cite{rosa2020translation}) and/or eye tracking (e.g. \cite{cui2021consultation}), in studies conducted in natural settings they tend to be used individually (e.g. \cite{kujamaki2010bildschirmvideos}) or in combination with retrospective interviews (e.g. \cite{albl2017professional}). 

Observation protocols were produced by combining non-participant observation and screen recordings. This is a research method that allows the researcher to gather data on a phenomenon while observing activities in their natural context without taking part in the observed activities \citep[609]{liu2010nonparticipant}. Protocols were used to track and reproduce the screen activity of the participants to shed light on the translation and PE process of GFL. \textcite[536]{munoz2019translating} maintain that broken words, pauses, changes, and searches may generally be indicators of translation problems. In the present work, they were used as types of screen activity to study the translation and PE process of participants. Broken words refer to translators leaving words unfinished. Changes are text modifications which correspond to a change in the original planning of the translation, which does not include typos.

Retrospective interviews were used to gather data on the perceived difficulty as well as the personal experience of participants. The use of TAP and/or retrospective interviews originally comes from psychology \citep{ericsson1984protocol} and was first applied to TS by \textcite{krings1986kopfen}. Some concerns arose about the validity of the method. These include the completeness and reliability of verbal reports and their possible effect on cognitive processes. They also cannot yield information about unconscious processes (see \cite{jakobsen2017translation}). Nevertheless, TAP and retrospective interviews are used in TPR to gain further insights into translation tasks \citep{dimitrova2009exploring}. They are of great value, especially if used to triangulate data \citep{jakobsen2017translation}, as in the current experiment with screen recordings and observation protocols. 

In addition, the produced target texts were annotated. Since the focus of this study is not on translation quality but rather on the ease of integrating GFL into the translation and PE process, only the selected GFL strategies and the success of their use were annotated. Further details on the study preparation, data collection and analysis are described in the next section.

%----------------------------------------------------------------------------------------
%	SECTION 2
%----------------------------------------------------------------------------------------
%qui scelte specifiche in devising the method and study
\section{Study preparation}
In the following subsections, different aspects of the study preparation are described in detail. These include criteria for language, text, assignment, and participant selection.

\subsection{Languages and GFL strategy} \label{sec:texts}

English was selected as the source language because it is a notional gender language that carries inflection only in third-person singular pronouns and few nouns. Reference to all genders is hence relatively straightforward and GFL strategies are more common. Notably, singular \textit{they} has nearly become the standard for referring to non-binary individuals \citep{gendercensus2021} and has historically been used to refer to people of unknown gender \citep{baron2020s}. Also, in English-speaking countries, especially the USA, UK, Canada, and Australia, awareness of non-binary genders seems to be more widespread with, for instance, celebrities coming out as non-binary and TV series portraying non-binary people. 

German was selected as the target language since it has grammatical gender and requires extensive gender marking in articles, nouns, and adjectives. Furthermore, a multitude of GFL strategies exist in German and no standard has yet been widely adopted. Due to this multiplicity of GFL strategies as well as differences in language structures, achieving linguistic inclusion or neutralisation of all genders is challenging for both human and MT from English into German. The language pair English-German is also a good selection for the MT-based PE component of this study because MT is known to work better for high-resource languages due to the amount of training materials available \citep{forcada2017making}.

As shown in Figure \ref{fig:coding}, GFL strategies to be used in the translation and PE experiment included (i) gender-neutral rewording, (ii) gender-inclusive characters, and (iii) neosystems. In the figure, different word classes are highlighted in different colours: violet for adjectives, yellow for singular non-binary nouns, green for plural nouns, light blue for articles, and red for pronouns. Each GFL approach differs in its application and was therefore expected to present different challenges. While (i) and (ii) are already used to a certain extent, especially to avoid masculine generics, (iii) represents a new approach to gender fairness and is devised by and for non-binary people. Three texts were then retrieved based on several selection criteria, i.e. text typology, variety of gendered elements, length, and readability.

\begin{figure}[ht]
    \centering
    \includegraphics[width=0.8\textwidth]{figures/coding.pdf}
    \caption{Extracts from the translations according to the three different study conditions}
    \label{fig:coding}
\end{figure}

\subsection{Texts and assignments} \label{sec:assignments}
Since non-binary individuals are increasingly featured in TV series and films, they also appear in popular texts, such as magazine articles \citep{misiek2020misgendered}. General language texts on non-binary actors starring in English-language TV series were retrieved online. Another reason for choosing this type of text was that they present and/or describe non-binary individuals in positive terms. The first text is about the Netflix series \textit{Sex Education}, where non-binary actor and singer Dua Saleh played the non-binary student Cal. The source of the text is the online website of the American magazine \textit{Newsweek}.\footnote{\url{https://www.newsweek.com/finally-sex-education-season-3-praised-groundbreaking-non-binary-representation-1629326}} The second text is about the promotion of non-binary actor E.R. Fightmaster as a recurring cast member in \textit{Grey's Anatomy}. It has two sources, namely the online British news site \textit{Metro}\footnote{\url{https://metro.co.uk/2021/10/27/greys-anatomy-casts-e-r-fightmaster-as-first-non-binary-doctor-15498268/}} and the movie and TV news website \textit{Screen Rant}.\footnote{\url{https://screenrant.com/best-non-binary-gender-nonconforming-characters-in-television/}} The third and final text presents the Canadian TV series \textit{Sort Of}, both created and interpreted by non-binary artists. Also in this case, two sources were used to create the text, i.e. the Canadian pop culture website \textit{Complex}\footnote{\url{https://www.complex.com/pop-culture/bilal-baig-cbc-sort-of-interview}} and the Canadian news site \textit{CBC}.\footnote{\url{https://www.cbc.ca/arts/new-show-sort-of-doesn-t-just-break-representational-barriers-it-s-also-more-than-sort-of-great-1.6186771}}

%In addition, research found there are no significant differences in time between translation and PE (see for example \cite{jia2019does}). Temporal effort could be then used to compare the impact of GFL on translation and PE.

Texts were also selected based on how many and which kind of personal references to non-binary people and mixed-gender groups they contained. Each text contains singular \textit{they} and at least one inflected form, two nouns in the singular form as well as articles or relative pronouns in reference to a non-binary individual. In addition, each text also contains a noun in the plural form referring to a mixed-gender group or non-binary individuals. Texts were also slightly modified to comply with these requirements. Study materials are included in the Appendix. Here is an extract:

\begin{quote}
Leading the way is \textit{a non-binary actor and musician} (non-binary article, adjective, and noun) Dua Saleh as Cal, \textit{a new student} (non-binary article, adjective, and noun) at Moordale (…) Hope is determined to turn Moordale's reputation as a ``sex school'' around and in doing so, she sets out to squash \textit{students'} sexual freedoms and \textit{their} right to express \textit{their} (mixed-gender group) identity. Cal stands up against Hope (…) \textit{They} (non-binary pronoun) refuse to wear the girls uniform and is reluctant to wear the tight-fighting boys uniform.\footnote{Grammar and spelling errors in the original are reproduced here.}
\end{quote}

In order to annotate the target texts, the gender phrases they contained (nine, twelve, and ten respectively) and their translations were identified and analysed.

In TPR involving the use of eye tracking, short texts of around 200--300 words are usually selected, for example, because the researcher must analyse a significant amount of data and participants’ motivation may decrease because of tiredness or boredom \citep[261]{o2009eye}. This also applies to the present study as participants were required to translate or post-edit three different texts, one after the other. As a consequence, only the first paragraphs of the retrieved texts were used in the study. They consist respectively of 154, 149, and 152 words. Details on the texts are summarised in Table \ref{tab:texts}.

\begin{table*}[t]
    \caption{Details on translation and PE materials}

    \resizebox{\textwidth}{!}{\begin{tabular}{ccccc}
    \lsptoprule
       \textbf{Text No.} & \textbf{TV Series} & \textbf{Instructed Gender-Fair Strategy} & \textbf{Word Count} & \textbf{Gender Phrases} \\
       \midrule
        1 & Sex Education & Gender-Neutral Rewording & 152 & 9 \\
        2 & Grey's Anatomy & Gender-Inclusive Characters & 151 & 12 \\
        3 & Sort Of & Neosystems & 163 & 10 \\
        \lspbottomrule
    \end{tabular}}

    \label{tab:texts}
\end{table*}

To ensure comparability of estimated translation and PE times, readability scores were computed using the Flesch-Kincaid readability test \citep{kincaid1975derivation} as in \textcite{lopez2016can}.  All texts were scored as ``fairly difficult to read''. This test accounts for the number of words and their length but ignores semantic difficulty. For this reason, the perceived difficulty of the source texts was discussed in the retrospective interviews. 

%Participants were instructed to select a different approach for each in order to test as many approaches to GFL as possible. As shown in Figure \ref{fig:translations}, these included: (i) gender-neutral rewording, (ii) gender-inclusive systems, (iii) neosystems (in the figure, different word classes are highlighted in different colors). Each approach differs in its application and was therefore expected to present different challenges.

Participants received a translation brief together with instructions on tasks, PE guidelines and a handout on GFL. This was designed to resemble a realistic order to aim for an authentic experimental setting (see materials linked in the Appendix). Therefore, the assignment description included client specifications, target audience, instructions concerning resources that could be used, and for the PE group also guidelines on full PE.

Austrian TV magazine \textit{TV-Media} was selected as the client for whom the translations were to be prepared. The translation brief specified that the source texts were general-language texts, thus prepared terminological resources would not be provided. The target audience was said to be general. Translators received the texts and had to deliver them as text documents in Word format. They were allowed to use their usual translation environment and tools as well as online resources. The texts were to be delivered as completed translations that would not need proofreading. 

Finally, people receiving the translation assignment were asked not to use MT. People who received the PE assignment were also given rules for full PE. The translation brief also included instructions on which approach to GFL was to be used. Five minutes before the video conference, participants received the texts with a shortened version of the translation brief and the texts to translate and post-edit.

%----------------------------------------------------------------------------------------
% SECTION 4
%----------------------------------------------------------------------------------------
\subsection{Participant selection} \label{sec:participants}
The number of participants, namely twelve, was established based on existing literature in the field of TPR (see for example \cite[255]{o2009eye}). O'Brien reviews previous studies in the field and notes how these feature a number of participants ranging from five to twenty. This is generally due to funding constraints. She also underlines how results from studies with twelve participants cannot be used to make generalisations, but have nevertheless significant value. Criteria for participant selection in the translation experiment were then established and aimed at ideal conditions in the translation market.

Translators needed (i) to have at least two years of work experience in the translation industry; (ii) to be highly proficient in English (C1+ level of the Common European Frame of Reference for Languages (CEFR)), i.e. the source language; (iii) to be German native speakers, as the translation into the L1 is most common; and (iv) to hold a relevant degree, e.g. translation/interpreting, linguistics, languages. An additional criterion was limited knowledge of GFL.

Previous joint studies on this topic \citep{burtscher-2022-genderfairmt,gromann2023participatory} suggest that few professional translators in Austria have extensive experience with GFL beyond the binary. As this study strives for representativity, translators with little or no knowledge of GFL were sought. This criterion could only be met in part. First, it does not apply to non-binary individuals who, for obvious reasons, were expected to be more sensitive to the topic and have higher expertise. Second, the topic is particularly sensitive and still controversial: only people with positive attitudes towards or a general interest in GFL, i.e. with reasonable expertise in the topic, participated in the study. %The recruitment of participants was more difficult than expected and most of the people willing to participate had, if not positive feeling towards it, at least a certain interest for GFL use and translation. 

Participants were selected by means of opportunity sampling. A call for participation in the translation experiment (see Appendix) was disseminated on the main social media platforms, i.e. Facebook, Twitter, Instragram, and LinkedIn, used by the author and by UNIVERSITAS. UNIVERSITAS also disseminated the request in their mailing lists and newsletters. The association was involved since it supported the study financially. Candidates were paid and worked from their familiar environment to increase the experimental realism of the study. The call for participation was disseminated from 31 May 2022 onward. Originally, a deadline for the application had been set for 6 June 2022. Since few translators had responded, the deadline was extended to 21 June 2022. After that, the author started to search for participants in his and his co-supervisor's personal network. 

An additional selection criterion was gender. Since translators are mostly women, eight out of twelve participants in this sample were women. Three men and one non-binary participant were selected in order to create a sample that was as representative as possible. However, no statistics concerning translation professionals in Austria were available, thus limiting the representativity of the study. To collect information concerning the participants' profiles, a pre-study questionnaire was created with Google Forms (see Figure \ref{fig:questionnaire}), consisting of twenty open and closed questions on demographics, work, GFL and PE experience as well as GFL use. The profile of the candidates is further discussed in Chapter \ref{results}. 
\begin{figure}[ht]
    \centering
    \includegraphics[width=0.7\textwidth]{figures/questionnaire.pdf}
    \caption{Screenshot of the questionnaire for participants}
    \label{fig:questionnaire}
\end{figure}

%----------------------------------------------------------------------------------------
%	SECTION 5
%----------------------------------------------------------------------------------------
\section{Data collection and analysis}
In the following subsections, the data collection and analysis process is described in more detail, delving into interview preparation, transcription and analysis, translation times measurements, observation protocols preparation and analysis, and segment-based target text annotation. 

\subsection{Interviews}
Semi-structured interviews were conducted immediately after the study. This format was selected as it allows for the author's flexibility to study specific aspects of the participants' experience in depth. The interview guide consisted of twelve questions for the group of translators and fourteen questions for the group of post-editors, respectively. These included impressions of the tasks, such as the challenges faced, the participants' preferred GFL approach, the impact of GFL use on their productivity during the translation and PE process, as well as their use of MT in general for translators, and within this study for post-editors (see Appendix for the specific questions). Specifically, post-editors were asked how their experience with PE was, whether the MT draft made their work easier and/or faster, and whether they would consider using MT again in the future for a similar assignment. 

In studies involving the use of eye tracking and keylogging, the log file generated during the translation experiment is at times replayed and the writing process is disclosed to the participants \citep[112]{dimitrova2009exploring}. In this study, the target and source texts were used as cues, therefore participants were allowed to read them during the interviews. Interviews lasted approximately 45--50 minutes, with deviations among participants. 

Data were collected from July to October 2022. At the same time, the interviews were transcribed following the rule system proposed by \textcite{dresing2008qualitative}, which was later updated by \textcite{dresing2012praxisbuch}, who created a content-semantic transcription system that prefers readability over details concerning pronunciation, colloquialisms, and tone. This decision was taken since elements such as pronunciation and the use of colloquialisms are not relevant to the scope of the present study. The transcripts of the interviews were also sent to participants to comply with good scientific practice.\footnote{In compliance with regulations at the University of Graz, ethical approval was sought and received for the present study (Ethikantrag Nr. 87 - 2021/22).}

The interviews were then analysed in MAXQDA,\footnote{\url{https://www.maxqda.com}} a software for computer-assisted qualitative and mixed-method data analysis, and following the basic process of thematic content analysis by \textcite{kuckartz2014qualitative} as depicted in Figure \ref{fig:qualitativecontentanalysis}. During the analysis, suggestions and good practices for work in MAXQDA proposed by \textcite{kuckartz2020fokussierte} were also followed to ensure a systematic analysis of qualitative data. For instance, a project memo with the research topic, research questions, and information related to the study participants was created. Also, the project logbook was used to track the progress of the data analysis. During the first reading of the interviews, memos were added and text passages were highlighted with colours or with emojis, as shown in Figure \ref{fig:maxqda1}. The transcripts were then summarised. 

\begin{figure}[ht]
    \centering
    \includegraphics[width=0.8\textwidth]{figures/qualitativecontentanalysis.pdf}
    \caption[Steps for thematic qualitative text analysis]{Steps for thematic qualitative text analysis as in \textcite[70]{kuckartz2014qualitative}}
    \label{fig:qualitativecontentanalysis}
\end{figure}


\begin{figure}[ht]
    \centering
    \includegraphics[width=0.8\textwidth]{figures/maxqda1.pdf}
    \caption{First text reading in MAXQDA}
    \label{fig:maxqda1}
\end{figure}

After this step, a first category system was inductively developed based on the interview guide. For almost every question, a category was created and defined in the code memos. Additionally, three other categories were developed following \textcite{kuckartz2020fokussierte}, i.e. \textit{other}, \textit{interesting}, and \textit{quotable passages}. The \textit{other} category was used for aspects that are relevant to the research questions but did not belong to a category yet. The \textit{interesting} category was used for passages that are not fully relevant to the research questions but still cover potentially interesting aspects for the present work. The \textit{quotable passages} category was used to select statements to be quoted in the result section of the work. Basic coding was then performed on all interviews. 


In the third step, segments assigned to each category were read to expand the category system with subcategories. During the fine coding phase, an order of category relevance was followed to create subcategories and then further develop and/or adapt them to better fit the data. The coded segments were consequently read numerous times until a saturation point was reached, i.e. no further subcategories were added anymore. The final category system comprises eleven main categories: 
\begin{enumerate}
\setlength\itemsep{-0.5em}
    \item \textbf{ST -- Challenges in Source Texts}: used when participants talk about the source texts in general and not about translation-related problems.
    \item \textbf{GNL -- Gender-Neutral Language}: used when the translation process of the first assignment, i.e. the translation with gender-neutral rewording, was discussed.
    \item \textbf{GIC -- Gender-Inclusive Characters}: used when participants elaborated on the translation process of the second assignment, i.e. the translation with gender-inclusive characters.
    \item \textbf{N -- Neosystems}: used for impressions concerning the translation process during the third assignment, i.e. the translation with neosystems.
    \item \textbf{GI -- General Impressions}: used when participants discussed personal impressions of the study such as what they liked doing or not.
    \item \textbf{MT -- Use of MT and Dis/advantages}: used when participants elaborated on the use of MT in general as well as for texts featuring GFL.
    \item \textbf{PEMT -- Impressions of the PE Task}: this category was applied to interviews with translators who were assigned the PE task and therefore discussed their impressions regarding their speed, the quality of the MT outputs and similar matters.
    \item \textbf{SP -- Speed and Productivity}: used for the participants' impressions of the impact that the use of GFL exerted on their work speed and productivity.
    \item \textbf{O -- Impact of Observation}: used when participants reflected on how they behaved during the study because of the researcher's observation.
    \item \textbf{RA -- Real Assignments}: used for discussions on how the participants would approach the translation or PE of GFL in a real-world scenario.
    \item \textbf{D -- Desires}: used when participants expressed their desires for optimal work with GFL.
\end{enumerate}
Finally, MAXQDA was used to compare the coded data and to create tables that will be presented in Chapter \ref{results}. 

%----------------------------------------------------------------------------------------
%	SECTION 6
%----------------------------------------------------------------------------------------
\subsection{Segment-based target text annotation}  \label{sec:annotation}
After data collection, target texts produced by the participants were annotated for GFL strategies used and the success of their application. The devised method was inspired by approaches to both TQA (\cite{castilho2018approaches}) and qualitative content analysis (\cite{kuckartz2014qualitative}). The process is explained in more detail below. 

First, two Excel documents were created, one for translation and the other for PE analysis. For each source text, nine, twelve, and ten gender phrases were identified with references to groups, e.g. ``audiences'', and non-binary actors/characters, e.g. ``non-binary neuroscientist'', that required gender marking in German. These were pasted into the Excel sheets along with all corresponding translations. Following typical approaches to qualitative content analysis \citep{kuckartz2014qualitative}, the texts were read several times and a basic category system was developed. The category system differed for each text since participants were instructed to use a different gender-fair strategy, namely gender-neutral rewording, gender-inclusive characters, and neosystems. 

The system included different GFL strategies, instances of misgendering, use of masculine generics, and mistakes in the application of GFL. Each category was assigned a number, as shown in Table \ref{tab:strategyannotation}. Texts were accordingly annotated, as in Figure \ref{fig:translationannotation} and the category system was refined. During this step, some categories were merged to provide a clearer overview of the results. 

Changes included merging the categories \textit{rewording}, \textit{abstraction}, and \textit{change of perspective} into the general category \textit{rewording} for the first text, and the use of the broad category \textit{other strategy} for gender references in the second and third text for which participants did not use gender-inclusive characters and neosystems. For this reason, in Table \ref{tab:strategyannotation} some numbers and categories were removed. Finally, a count function in Excel was used to calculate the frequency of each category both in general and for each analysed gender phrase. 

\begin{table}[ht]
    \centering
    \caption[Target text annotation categories]{Numbers assigned to categories for the translation annotation of the first text}
    \begin{tabular}{cc}
    \lsptoprule
       \textbf{Code}  & \textbf{Strategy}  \\
       \midrule
        0 & No strategy \\ %\hline
        1 & Indefinite Pronoun \\ %\hline
        2 & Gender-Neutral Words \\ %\hline
        4 & Rewording \\ %\hline
        5 & Collective Noun \\ %\hline
        6 & Participial Form \\ %\hline
        7 & Name Repetition \\ %\hline
        10 & Pronoun Omission \\ %\hline
        11 & Misgendering/Mistranslation of \textit{they} \\ %\hline
        12 & False Term \\ %\hline
        13 & Gender-Inclusive Characters \\ %\hline
        14 & Relative Pronoun \\ %\hline
        15 & Masculine Generics \\ %\hline
        16 & Omission \\
        \lspbottomrule
    \end{tabular}
    \label{tab:strategyannotation}
\end{table}

\begin{figure}[!ht]
    \centering
    \includegraphics[width=0.8\textwidth]{figures/translationannotation.pdf}
    \caption{Example of target text annotation}
    \label{fig:translationannotation}
\end{figure}

Similarly, the machine translations produced with DeepL were annotated for bias. For each text, the nine, ten, and twelve phrases containing gender references to be post-edited were selected and pasted into the previously mentioned Excel sheet. The MT outputs were annotated for three types of bias: (i) misgendering (i.e. a binary gender was assigned to a non-binary person mentioned in the text), (ii) mistranslations of singular \textit{they} (i.e. the pronoun use was misinterpreted and translated in the third-person plural into German), and (iii) masculine generics. 

%----------------------------------------------------------------------------------------
%	SECTION 7
%----------------------------------------------------------------------------------------
\subsection{Translation times and observational protocols}
%In Chapter 4.2. and its subchapters to me the actual translation/PE process is missing - maybe add one chapter on how it was planned/designed? How much time/When obtain the text, video recording, etc.? Maybe combined with 4.2.6?

Screen recordings are considered an appropriate method to elicit data on the translation process, even though they are not as precise as keylogging. They allow the researcher to collect additional data on the type of consultations made by translators and to see what happens on-screen. For this reason, they are regarded as a valid method for data triangulation \citep{angelone2013impact}. In the present study, they were used to measure translation times as well as to complete observation protocols produced with the researcher's notes.

Translation times were measured and imported into an Excel sheet, together with the demographics and data on the use and perception of GFL by participants as collected in the pre-study questionnaire. The quantitative analyses in this study aimed to examine whether different tasks (translation vs. PE) and GFL strategies (gender-neutral rewording, gender-inclusive characters, and neosystems) had an impact on participants' performance. To achieve this, both descriptive and inferential statistical methods were employed, using Microsoft Excel for basic analysis and the statistical analysis software Jamovi\footnote{\url{https://www.jamovi.org}} for more advanced statistical testing.

Descriptive statistics, including the mean (average), median (middle value), and standard deviation (a measure of how spread out the values are), were calculated to summarise the translation and PE times. These metrics help describe the general tendencies in the data and provide a foundation for further statistical testing. This approach is in line with standard practice in TPR \citep{o2009eye}.

Inferential statistics are used to determine whether observed differences in the data are statistically significant, that is, unlikely to have occurred by chance alone. Before selecting appropriate tests, assumptions such as normal distribution (whether the data follow a bell-shaped curve) and homogeneity of variance (whether different groups have similar variability) were checked. Since these assumptions were not met, alternative robust tests were applied.

Inferential statistics were performed importing the Excel sheet with the experiment data into Jamovi. The different types of variables, e.g. age, gender, and translation times, were assigned to different categories, i.e. nominal, ordinal, and continuous. 

A \textit{Welch's t-test} was conducted to establish whether the task, i.e. translation and PE, affected the participants' speed. Unlike the standard t-test, Welch’s test does not assume equal variances between groups, making it more appropriate for small or unequal samples \citep{ruxton2006}. Because the data were also not normally distributed, a \textit{rank transformation} was applied. This means that the actual data values were replaced by their rank positions (e.g. 1st, 2nd, 3rd), which reduces the impact of outliers and skewed data while allowing parametric test structures to be retained \citep{zimmerman1993}. This method is recommended in TS when working with small and heterogeneous samples \citep{mellinger2016quantitative}.

To examine whether the different GFL strategies had an effect on processing times, a \textit{repeated measures ANOVA (rANOVA)} was used. This test is suitable when the same participants complete multiple tasks, allowing within-subject comparisons. However, rANOVA assumes sphericity (that the variances of the differences between all pairs of conditions are equal). Since this assumption was violated, the \textit{Greenhouse-Geisser correction} was applied to adjust the degrees of freedom and reduce the risk of false statistical significance \citep{girden1992}. This procedure is common in psycholinguistic and TPR experiments involving multiple within-subject measurements.

Following a significant rANOVA result, \textit{post-hoc tests} were conducted using the \textit{Bonferroni correction} to adjust for multiple comparisons. When several comparisons are made, the chance of finding a significant result by chance increases. The Bonferroni method addresses this by lowering the threshold for significance, thus preserving the overall error rate \citep{bland1995}. Although conservative, this method ensures the reliability of the results and is recommended for studies with multiple conditions.

Screen recordings along with observation notes were used to produce observation protocols. The source texts were split into segments based on their punctuation and entered in an Excel sheet. The aim of the protocol analysis was to summarise the translation and PE process, more specifically the activities needed to integrate GFL in the target texts. \textcite[536f]{munoz2019translating} consider broken words, changes, searches, and pauses as indicators for possible translation problems and hence cognitive processes. Such activity categories were used in the analysis of the observation protocols.

While in \textcite{munoz2019translating} the first category was \textit{broken words}, for the present work, noun phrases, e.g. ``non-binary actor'', were regarded as larger units that are translated together and hence the category \textit{broken units} was preferred. The annotation included frequencies of each phenomenon as well as explanations of what happened on the screen, e.g. which sources were searched for and when pauses were taken during the translation and PE process, as shown in Figure \ref{fig:observation}. A column was also added for general comments and observations that could potentially provide more insights into the participants' workflows. 

\begin{figure}[ht]
    \centering
    \includegraphics[width=0.8\textwidth]{figures/observation.pdf}
    \caption{Example of observational protocols}
    \label{fig:observation}
\end{figure}

In TPR, pauses have been operationalised differently (see \cite{kumpulainen2015operationalisation}) based on the methodology used in studies as well as on their objectives. In screen recordings, it is almost impossible to measure pauses of less than three seconds. Hence, only pauses of three seconds or more (3+) were counted for the analysis. Finally, absolute and relative frequencies of the analysed activities were computed and compared among participants, assignments, and text segments. The absolute frequencies of activities per participant were also pasted into the aforementioned Excel sheet containing the data of the translation and PE times. 

In this case as well, inferential statistics were performed. As in the analysis of the translation times, two tests were conducted. The first was a Welch's test on rank-transformed data to test whether the task type affected the number of activities observed in the screen recordings. The second was a rANOVA with Greenhouse-Geisser correction and a post-hoc test with Bonferroni correction to test whether the assignments affected the number of the activities tracked.

This chapter outlined the methodology of the present study, which integrates both quantitative and qualitative methods from process- and product-oriented research in TS. Twelve participants were recruited and divided into two groups: one was tasked with translating, and the other with post-editing three texts related to TV series featuring non-binary characters. The translation and PE process was recorded, and the author conducted observations while taking detailed notes. Processing times were measured, and participants were interviewed after each study session. Finally, the target texts were analysed with particular emphasis on the strategies employed and the errors made.
